Een gefrustreerde klier
plaste eens door een kier

van een houten voordeur.
Zijn baldadig humeur

kwam voort uit zoete wraak
voor de klap op zijn kaak

vanuit Emmy’s vader.
Hij werd des te kwader,

toen vader zijn ouders,
als zijn toezichthouders,

aansprak op hun zorgplicht.
Het verlies van gezicht

werd flink gecompenseerd,
want hij had hen onteerd

door zich af te trekken
achter schapenhekken.

Emmy die ging plassen
vlak achter de kassen

was zijn erotisch beeld.
Het had weinig gescheeld

of ook hij was ontsnapt.
Als enige betrapt

van zijn kameraden,
mocht hij zich ontladen

tegen het vermanen.
Een van zijn kompanen

klopte toen de klopper.
Ondanks zacht gemopper

kon de schelm zijn plassen
niet ineens aanpassen.

Vader ontsloot de deur,
terwijl de pislikeur

over zijn sokken liep
en de snaak om hulp riep.

Spreuken 26
vervloekte hem bitsig.

Via "Wie een kuil graaft”,
werd de revanche gestaafd.

De zaal was vierkantig.
De sfeer was parmantig.

Hij lag plat op zijn buik.
Zij depte in een kruik,

met haar onderwerping
als natte spiegeling,

eerst haar zachte handen,
alvorens die landden

op zijn behaarde rug.
Hij trok zich kort terug,

omdat de aanraking
kwam als een verrassing.

Haar makke massage,
de blote blamage

liet de landheer keren
tijdens het masseren.

Rondom stonden vazen
als volle wijnglazen.

Gluurders wilden weten,
onderuit gezeten

in hun chaise longues,
hoe hij zijn homspuitbus

in haar moeder bonkte.
Hoe de landheer pronkte

met zijn nieuwste jachtbuit.
Die zag er prachtig uit.

Het was sociaal leren
door te observeren,

internaliseren
en te imiteren.

Zie het theorema
van Albert Bandura.

Emmy las obsessief
haar zusters eerste brief.

“Ook in mama’s geval
stonden daar aan de wal

weer de beste stuurlui.
Er barstte een pensbui

los na slechts elf stoten.
Eenmaal volgespoten

juichten de getuigen.
Ze moest hem schoonzuigen

als slaafse afronding
van mama’s ontmaagding.

Nadat hij haar besprong
bleek zij en de nacht jong.

De kijkers dropen af,
Mama had geen vrijaf.

Het kraaien van de haan,
liet haar van hem afgaan.”

Emma vergeleek breed,
wat de landheer ooit deed,

als vuige profiteur
van zijn droit de seigneur,

met de situatie
tijdens de receptie

direct na de oproep.
Het was dezelfde groep

van spiedende gasten.
Hij liet haar betasten

tijdens de reünie,
vlak na zijn eruptie,

in plaats van dat de kliek
als werkeloos publiek

meteen moest verkassen.
Ook zij mochten plassen

hun opgespaarde zaad
in moeders dunne naad.

Want na twintig jaren
mochten ze gaan paren.

“Mama werd daar een hoer.
Zijn bordeel haar werkvloer.”

Zo raakten in die tijd
velen hun vroomheid kwijt,

niet door hun slechte aard,
maar door wat hen bezwaard.

Toch zag Jean-Jacques Rousseau
het ambitieniveau

van de mens als zuiver.
Zijn waarlijke huiver

kwam van de maatschappij,
die in zijn heerschappij

de mens steeds corrumpeert
en het kwaad stimuleert.

Haar tante gaf Emmy
de correspondentie

die vader behouden
achter had gehouden.
De brieven stonden vol
met Emma in de rol

van een prostituee.
Emma nam Emmy mee

in haar leven als hoer.
Met kronkels als richtsnoer

en grillen als inhoud
koos hij voor zelfbehoud.

Al kon hij niet lezen,
toch deed vader vrezen

dat de vele brieven
Emmy zou gaan grieven.

Emma had het inzicht,
waarom een mens vaak zwicht

voor seksuele kicks.
Resulterend in tics.

Ze zag het sublieme
vaak in het perverse.

Jacques Lacan bij uitstek,
maar ook Slavoj Žižek

zou haar zien als subject
dat optreedt als object

van iemands verlangen,
zonder haar te vangen.

Ze gaf zich vrijwillig.
Ze deed het gewillig.

Ze had snel opgepikt,
waarvoor elke klant knikt

en zo zichzelf ontdekt.
Ze kon dit intellect
echter niet spiegelen.
Dan werd het giechelen.

Ze had geen voorkeuren.
Alles mocht gebeuren.

Want Emma’s omgangsvorm
had geen vanillanorm.

Ze zag een fantasie
als helende fictie,

niet als wat men echt wil.
Sex was haar wonderpil

en haar erotisch spel,
dat van een seksrebel.

De roerige Emma
schreef over haar Lola

in de geest van de Kinks,
ook suggestief maar slinks.

Des landheers gemalin,
werd haar beste vriendin

en haar beste minnaar.
Was dit allemaal waar,

vroeg Emmy zich vaak af?
Emma verhaalde straf

en verzon vaak feiten.
Waarom zou ze pleiten

de landheer te mijden
wegens fysiek lijden?

“Taboes behelzen niets.
Iedereen heeft wel iets

dat specifiek opwindt
wat een ander raar vindt.

In deze voorkeuren
kan men pijn bespeuren

en door te verwennen
beter leren kennen.

Ook mama en papa
hebben elk hun trauma

dat zich uit in genot!
Ik weet zeker dat God

ons seks heeft gegeven
om ons te vergeven.

Ons daar verlossen gaat
van het gedane kwaad.

Ook jij zult ontdekken
jouw duistere plekken

en die van anderen.
Laat je veranderen

en vraag jezelf eens af,
wat stoot jou aan of af?

Ervaar jij opwinding?
Heb jij al ervaring?

Heb jij een fantasie?
Wees eens eerlijk Emmy!”

Emma - als Coyote,
de Navajo mythe -

overtrad de regels.
Hoorde bij de vlegels,

die zich nimmer schamen
en niet gehoorzamen.
Emmy las dit schrijven
tijdens het verblijven

bij haar oom en tante.
Vooral de pikante,

verborgen motieven,
in haar zusters brieven,

deden haar afvragen
of ze het kon wagen

tante in te lichten
over de berichten.

Zij waren ook mensen
met intieme wensen.

Door welke fetisjen
ontdekte Emmy hun

onverwerkte wonden?
Hun verborgen zonden?
